\documentclass{article}
\usepackage{amsmath,amsthm}
\newtheorem{theorem}{Theorem}
\begin{document}
\title{When recovery traffic sustains a settling fault}
\maketitle

\paragraph{Interface.}
A shared fixture starts degraded. In a slot containing an intrusive diagnostic,
its quiet-age is reset to zero. In a slot without one, the age grows by one.
The first attainment of the threshold $q\ge1$ restores joint service. The
mechanical analogy is a shared test contact that must be released while it
settles. This is an authored mechanism, not evidence of a particular robot or
firmware fault. It models an observation action that changes the plant and
not merely a read of its state. In this model settled service latches; the rolling
observation law continues. If future tests could reopen settled service, the
theorem would only bound each reacquisition, not continuous uptime.

Let $H\ge 1$. There are $n$ intrusive diagnostic streams, each requiring a completed check
in every rolling $H$ slots. At most $c$ checks can use the intrusive fixture per
slot. Checks are instantaneous at the modeled slot cut; they reset before the
quiet update. An upgraded stream can instead use a distinct, nonintrusive
evidence path, with independent adequate capacity and the same freshness.
The theorem does not infer such a path from a status field. The initial
history is selectable by a periodic warm-up; an arbitrary preexisting set of
ages is a different problem. All integer-slot results use clean, reliable
operations and concurrent reservation of one shared cadence.

\begin{theorem}
For $0\le q\le H$, a quiet block of length $q$ is compatible with the
rolling checking law iff $n\le c(H-q)$. Hence the minimum number of streams
that must switch to the stated nonintrusive path is
$\max(0,n-c(H-q))$. For $q>H$, every stream must switch.
\end{theorem}
\begin{proof}
Extend the quiet block (or its first $H$ slots) to an $H$-slot window.
Every intrusive stream must execute at least once inside that window.
At most $H-q$ slots are usable if $q\le H$, so capacity gives the bound,
independently of the policy's memory or feedback. For $q>H$ the quiet
block contains an entire window. Conversely, in each $H$-slot frame place
the intrusive checks in the first available slots, up to $c$ per slot, each
stream at the same phase in each frame. Every $H$ consecutive slots contain
one occurrence of each phase. The remaining consecutive slots are quiet.
With zero intrusive streams, all slots are quiet. Substituting the retained
stream count yields the exact upgrade requirement.
\end{proof}

\paragraph{Feedback witness.}
Let $H=4$, $n=2$, $c=2$ and $q=3$. Two independent watchdogs probe on
timeout (every fourth slot while degraded) and occupy opposite phases.
Either watchdog in isolation allows a three-slot quiet interval. In the joint
execution there is a probe every second slot, so the fixture never leaves the state
that triggers probing. No continuing original disturbance is needed.
A shared cadence batches both probes in one slot out of four, keeps each check
fresh, and settles. A direct phase-aware scheduler with the same batching
capability predicts the same result. When $c=1$, shared cadence alone cannot
work; one stream must switch or the freshness or settling obligation must change.
When one purported read-only check actually resets the fixture, count that
stream as intrusive again; a last-check timestamp does not establish its effect.

\paragraph{Recovery from already running watchdogs.}
The periodic theorem does not grant a change to the pending deadlines. Fix
a snapshot with ages $a_j<H$ and next deadlines $d_j=H-a_j$ in slots numbered
from one. Restrict to a recovery cut $T=s+q\le H$, with the last $q$ slots
quiet. The evidence law is monitored through the cut; afterward a stipulated
adequate read-only mode serves every stream on time. Define $D_T=\{j:d_j\le T\}$.
Exactly the streams in $D_T$ need a check before the cut. There is a lawful
prefix iff
\[
       \big|\{j\in D_T:\min(s,d_j)\le t\}\big|\le c t
       \quad\text{for every }0\le t\le s.
\]
Any feasible schedule must meet these prefix counts. In nondecreasing order
of effective deadline, the $k$th required check can be put in slot
$\lceil k/c\rceil$; the inequalities give that slot by its deadline.
For $c=0$ the test permits only an empty due set. No second check of a
stream can be due by $T\le H$ after its first service in slot one or later.
For $H=4$, $c=2$, $q=3$, deadlines $(1,4)$ allow service in slot one
and settling at four; deadlines $(1,2,3)$ allow neither settling at three
nor at four despite aggregate service feasibility without quiet. Even two
streams with deadlines $(1,2)$ cannot settle at three, while deadlines
$(4,4)$ permit it. The carried deadline vector therefore matters even
when the periodic design has sufficient capacity. The test compares
against direct phase-aware search at identical information and costs.

\paragraph{Reuse and limits.}
The pigeonhole witness depends on which checks reset the same fixture, the
common rolling window, and their service capacity. Replacing a monitor while
preserving these interfaces preserves the bound; changing the checks'
availability phases or preexisting ages can invalidate the constructive cadence.
The capacity is a bound on distinct serviced streams, not on the maximum
observed concurrency in old logs. This is a finite scoped feedback model,
not a checked architecture disposition or field result.
\end{document}
